\section{Auditoría de fuentes e hilo conductor de la clase}

Esta clase abarca el lenguaje clínico, las secuencias de proteínas y la regulación génica. A primera vista parece un recorrido por cuatro artículos, pero el hilo conductor es muy unitario: \textbf{primero se escriben los objetos biomédicos como secuencias y después se adaptan las capacidades de representación, atención y generación del Transformer a las restricciones de cada dominio.} Por ello, esta clase no puede limitarse a extraer los resultados de Med-PaLM ni reducir la segunda mitad a «el Transformer también sirve para biología». Hay que preguntar en cada caso cuál es la entrada, cuál es la señal de supervisión, dónde están el contexto largo o el ruido, cómo se verifica la salida del modelo y qué no puede concluirse de los resultados experimentales.

\subsection{Cómo se reconstruye esta clase}

El sitio oficial del curso ofrece la grabación y tres trabajos recomendados, pero no publica un slide deck independiente. Esta clase toma como fuente visual principal la grabación oficial en alta definición completa de Stanford Online: el tamaño de imagen es de 1920 por 1080 píxeles; un muestreo cada dos segundos produjo 2,045 fotogramas, el filtrado por estabilidad y diferencia dejó 103 candidatos de alta exhaustividad y, tras compararlos uno por uno manualmente, se conservaron 84 estados didácticos distintos. Todas las imágenes se volvieron a extraer del video original a ancho completo, sin recortar el contenido del lado derecho; los subtítulos oficiales \texttt{en-US} hechos por personas se usan para recuperar las motivaciones, limitaciones, preguntas y respuestas en clase y transiciones del ponente.

\begin{longtable}{@{}L{0.18\textwidth}L{0.28\textwidth}L{0.42\textwidth}@{}}
\toprule
Capa de evidencia & Material local & Función en los apuntes \\
\midrule
Grabación oficial & Video de 1:08:09 y 84 imágenes didácticas & Determina el orden de la clase, las tablas, las fórmulas, las gráficas de resultados y la evidencia que el ponente mostró realmente \\
Subtítulos oficiales hechos por personas & \texttt{lecture17.en.srt} y timed transcript & Recuperan el «porqué», las correcciones de gráficas, la explicación de riesgos, el impacto de ingeniería y los juicios sobre el futuro \\
Artículos asignados & Med-PaLM, ProtNLM, Enformer & Contrastan la definición de las tareas, los mecanismos de los modelos, los objetos experimentales y los límites de las conclusiones \\
Artículos de contexto & Performer, DeepConsensus, PaLM, Flan-PaLM & Completan las fórmulas y el contexto de sistema que la clase muestra de forma condensada \\
Archivos de auditoría & manifest, coverage, ledger, blueprint & Demuestran que cada nodo visual y cada nodo de teacher voice tiene un destino explícito \\
\bottomrule
\end{longtable}

\begin{warningbox}{Límite histórico: 21 de febrero de 2023}
Esta clase reconstruye el estado del conocimiento del día en que se impartió. Med-PaLM 2, la serie Gemini, los modelos biomédicos multimodales, los eventos regulatorios y los benchmark más recientes, todos posteriores, no se incorporan retroactivamente como evidencia de la clase. Lo que el ponente plantea al final es una visión de investigación de foundation biomedical AI, no un sistema terminado que ya unifique las tareas clínicas, de proteínas y genómicas.
\end{warningbox}

\subsection{Un recorrido en tres etapas: clínica, proteínas, genoma}

La portada solo presenta el tema, pero la página de agenda revela el diseño de la clase: primero responde desde los primeros principios «por qué el Transformer», después entra sucesivamente en clinical, proteins y genomics, y al final discute el futuro. Al leer estas dos páginas conviene notar que el orden de los dominios desciende, en términos generales, por niveles de abstracción: del lenguaje natural que usan médicos y pacientes a las secuencias de aminoácidos, y de ahí a los nucleótidos y las señales regulatorias; pero en cada nivel que se desciende la salida del modelo debe redefinirse, y no pueden heredarse los criterios de evaluación de un chatbot.

\lecturefigure{slide-01-title-transformers-biomedicine.jpg}{Biomedical Transformers: título de la clase y ponente}{Diapositiva V001 recuperada de la grabación oficial; Stanford CS25 Lecture 17}

\lecturefigure{slide-02-agenda.jpg}{Ruta de la clase: aplicaciones clínicas, proteínas, genoma y futuro}{Diapositiva V002 recuperada de la grabación oficial; Stanford CS25 Lecture 17}

\teachervoice{Al iniciar, el ponente dice que su objetivo es «peel back the curtains», para que los estudiantes vean la innovación que está ocurriendo en la intersección entre la IA y la biomedicine. No pretende demostrar que un solo modelo domina todas las tareas, sino que elige casos que atraviesan el biomedical stack, de modo que los mecanismos comunes y las diferencias entre dominios se hagan visibles a la vez.}

\begin{importantbox}{Cinco preguntas que recorren toda la clase}
Para cada caso deben responderse en orden: ¿cómo se construye la secuencia de entrada? ¿Por qué es difícil el contexto? ¿De dónde proviene la señal de supervisión? ¿Cómo verifican la salida las personas o los experimentos? ¿Lo que respaldan los resultados es una mejora en un benchmark, un descubrimiento científico, un apoyo clínico o un despliegue real? Si estas cinco capas se mezclan, las puntuaciones de un modelo se presentan fácilmente, por error, como capacidad médica.
\end{importantbox}

\subsection{Resumen de la sección}

La cadena de evidencia de esta clase se compone de la grabación oficial, los subtítulos hechos por personas y cinco grupos de artículos principales. La estructura de la clase no es una «colección de chatbots médicos», sino una comparación de cómo el Transformer procesa distintas secuencias biomédicas, cómo cambia la complejidad, cómo define la salida y cómo los resultados del modelo se reconectan con la verificación clínica o experimental.
